% !TeX spellcheck = ru_RU
% !TEX root = vkr.tex
\section{Обзор}
\label{sec:relatedworks}

В разделе рассматриваются прототип модуля оценивания, методы его расширения
и готовые реализации. Сопоставление нужно для выбора состава пакета и
обоснования собственной разработки относительно готовых решений. Рассматривается точечное оценивание;
интервальные оценки и проверка гипотез в объём работы не входят.

\subsection{Прототип и исходная архитектура PySATL Core}
PySATL Core предоставляет параметрические семейства, несколько параметризаций
одного семейства и объекты конкретных распределений. Характеристики, например
плотность, функция распределения и моменты, могут вычисляться непосредственно
или через преобразования в графе характеристик \cite{reviewCore}.

Выбор PySATL Core как основы работы связан с возможностью исследовать
распределения и методы их оценивания при разных способах задания модели.
Граф позволяет получать недостающие характеристики из доступных:
например, вычислять функцию распределения интегрированием плотности.
Это сокращает число характеристик, которые нужно реализовать явно
при добавлении нового семейства. Несколько параметризаций позволяют
работать с одним семейством в разных координатах и сравнивать поведение
алгоритмов при их смене. Поэтому для выбранной задачи важно сохранить
эти возможности при подключении оценивателей, а не ограничиться
подгонкой распределений в заранее заданном представлении.

Оценивание должно использовать эту модель: принимать семейство и наблюдения,
определять свободные параметры и возвращать пригодное для дальнейших
вычислений распределение.

В курсовой работе разработан прототип максимального правдоподобия:
общий результат подгонки, проверка выборки, подготовка параметров и ограничений,
реестр формул, адаптеры оптимизаторов и диагностика
\cite{reviewPrototype}. Поддерживаются формульный и численный пути,
разные параметризации и частично фиксированные семейства. Отдельный прототип
эмпирических распределений предоставляет эмпирическую функцию распределения
(ECDF) и гауссову ядерную оценку плотности (KDE).

Главной проблемой, выявленной при реализации прототипа, оказалась работа
с фиксированными параметрами. В финальной версии планируется разработать
механизм управления ими, который позволит алгоритму оценивания работать
только со свободными параметрами, не зная, какие параметры зафиксированы
и каковы их значения. Подстановку фиксированных значений и проверку
связанных с ними ограничений должен выполнять этот механизм.

Эта основа решает задачу подключения отдельных способов оценивания к ядру,
но ещё не образует согласованный пакет всех выбранных методов. В частности,
моментные правила подготовки начальной точки MLE не являются самостоятельными
моментными оценивателями. Поэтому работа должна развивать
существующие прототипы, сохраняя подготовку задачи и численную инфраструктуру.

\subsection{Методы, планируемые к включению}
\label{sec:methods}
Ниже $x_1,\ldots,x_n$ обозначают независимые наблюдения,
$\theta$ --- свободные параметры модели.

\textbf{Минимизация критерия.}
\emph{Метод минимального контраста} минимизирует среднюю потерю
на отдельных наблюдениях:
\[
 C_n(\theta)=\frac1n\sum_{i=1}^n\rho(x_i,\theta).
\]
Квадратичная потеря приводит к методу наименьших квадратов и сильно
штрафует большие отклонения. Робастные потери, например потеря Хьюбера,
растут медленнее при больших отклонениях и позволяют уменьшить влияние
выбросов на оценку.

\emph{Метод максимального правдоподобия (MLE)} соответствует потере
$\rho(x,\theta)=-\log f_\theta(x)$, где $f_\theta$ --- плотность или функция
вероятностей. Минимизация $C_n$ равносильна максимизации правдоподобия.
Методу требуется вычисление $f_\theta$ и проверка принадлежности наблюдений
носителю модели. Планируется сохранить формульный и численный пути
из прототипа.

\emph{Метод минимального расстояния} сравнивает модель с выборкой целиком:
например, функцию распределения с эмпирической функцией распределения
(ECDF) либо вероятности дискретной модели с частотами. Выбор расстояния
и весов определяет, какие расхождения сильнее влияют на оценку
\cite{reviewDistrMod}.

\textbf{Согласование характеристик и оценочные уравнения.}
Моментные методы выбирают параметры так, чтобы характеристики модели
были равны вычисленным по выборке или отличались от них как можно меньше.

\emph{Обычные моменты} --- средние степеней случайной величины.
Первый определяет среднее, а второй вместе с первым --- дисперсию.
Для нормальной модели оценками служат среднее выборки и средний квадрат
отклонений от него. Моменты должны существовать; высокие порядки
чувствительны к выбросам.

\emph{Вероятностно-взвешенные моменты (PWM)} --- средние значений,
умноженных на веса, заданные через функцию распределения.
Например, возрастающие веса сильнее учитывают верхнюю часть распределения
\cite{reviewGreenwood}.

\emph{L-моменты} выражаются через PWM. Их выборочные оценки --- суммы
отсортированных наблюдений (порядковых статистик) с определёнными весами.
Первый L-момент равен
среднему, второй описывает разброс, высшие --- форму распределения.
Они менее чувствительны к крайним значениям, чем обычные моменты высоких
порядков \cite{reviewHosking}.

\emph{Усечённые L-моменты (TL-моменты)} обобщают L-моменты: при их определении
исключают заданные числа крайних порядковых статистик слева и справа.
Это ослабляет влияние хвостов; некоторые TL-моменты существуют даже
без конечного среднего \cite{reviewHoskingTL}.

При избытке условий их точное выполнение может оказаться невозможным.
Обобщённый метод моментов (GMM) минимизирует
$g_n(\theta)^{\mathsf T}Wg_n(\theta)$, где $g_n$ --- вектор выборочных
моментных условий, а $W$ --- весовая матрица \cite{reviewHansen}.
Так моментный подход связывается с минимизацией критерия.

Более общий путь --- решение оценочных уравнений
$\sum_i\psi(x_i,\theta)=0$. Для дифференцируемой потери $\psi$ может быть
её производной, но не каждое уравнение задаётся скалярным критерием.

Робастный пример --- M-оценка положения Хьюбера: ограниченный вклад
больших отклонений уменьшает влияние выбросов \cite{reviewHuber}.
Нужны явные порог и правило определения масштаба, а при решении
уравнений --- проверка невязки и допустимости корня.

\textbf{Сопряжённое байесовское обновление.}
Здесь модель дополняется априорным распределением параметров.
При сопряжённости апостериорное распределение принадлежит тому же
семейству, что и априорное, а его параметры обновляются через достаточные
статистики \cite{reviewDiaconis}. Планируются нормальное априорное
распределение для среднего нормальной модели с известной дисперсией
и гамма-распределение для интенсивности экспоненциальной модели.
Точечный результат --- выбранные среднее, медиана или мода
апостериорного распределения в заданных координатах.
Этот путь ограничен доступными сопряжёнными парами.

\textbf{Непараметрические оценки и повторное оценивание.}
Среднее, дисперсия, моменты и квантили оценивают характеристики
распределения непосредственно по выборке. Порядковые статистики
служат основой квантилей и L-статистик. ECDF и эмпирическая функция
вероятностей (PMF) задают распределение с массами в наблюдаемых значениях.
Ядерная оценка плотности (KDE) сглаживает эти массы; ширина окна определяет
степень сглаживания и существенно влияет на результат \cite{reviewScipyKDE}.

Бутстреп использует уже выбранный оцениватель: непараметрический вариант
создаёт выборки с возвращением из наблюдений, параметрический ---
из подобранной модели \cite{reviewEfron}. Планируется сохранять повторные
оценки и сведения о неудачах; построение интервалов не входит в работу.

\subsection{Готовые реализации}
\label{sec:implementations}
Рассмотренные средства решают задачу оценивания на разных уровнях:
SciPy предоставляет базовые вычисления, fitdistrplus и distrMod ---
согласованные средства подгонки в R, OpenTURNS --- развитую объектную
систему в Python, а statsmodels, Lmo и lmomco --- отдельные направления
оценивания. Сравнивать их только по наличию MLE недостаточно.

Для данной работы важны четыре свойства: охват выбранных методов;
возможность использовать одну модель при смене оценивателя;
работа с параметризациями и фиксированными параметрами;
согласованное представление результата и повторных вычислений.
Требуется, например, на одном частично заданном семействе сравнить MLE,
моментную и дистанционную оценки, а затем применить к каждой бутстреп.
При смене метода должны сохраняться ограничения модели и смысл
свободных параметров.

\subsubsection{SciPy}
Преимущество SciPy --- широкий набор распределений, готовые численные
процедуры и возможность добавлять собственные распределения через
\texttt{rv\_continuous} \cite{reviewScipyDistribution}.
В \texttt{rv\_continuous.fit} доступны MLE и обычный метод моментов,
фиксирование параметров и замена оптимизатора; \texttt{stats.fit}
возвращает диагностику и поддерживает MLE и максимум произведения
интервалов вероятностей \cite{reviewScipyContinuous,reviewScipyFit}.
Для стандартной подгонки этих средств часто достаточно.

Ограничение для рассматриваемой задачи --- различие контрактов:
\texttt{dist.fit} возвращает кортеж, \texttt{stats.fit} --- объект,
а GMM и TL-matching не входят в перечень методов этих двух операций.
Добавление нового распределения само по себе не объединяет эти способы
оценивания. Поэтому SciPy подходит как вычислительная основа, но
пользователю всё ещё приходится организовывать выбор метода,
преобразование параметров и сбор результатов.

\subsubsection{fitdistrplus и distrMod}
Преимущество \texttt{fitdistrplus} --- готовый процесс подгонки:
MLE, согласование моментов и квантилей, критерии согласия, бутстреп,
диагностика и графическое сравнение моделей \cite{reviewFitdistrArticle}.
Можно задавать пользовательские распределения и фиксированные параметры
\cite{reviewFitdistrManual}. Это близкий функциональный аналог,
а не просто набор формул. Ограничение --- его интерфейс рассчитан на R:
для использования в Python нужны выполнение R-кода и преобразование
данных и результатов. Готовую L-/TL-моментную ветвь в данном сравнении предоставляют
специализированные пакеты Lmo и lmomco.

\texttt{distrMod} ближе к требуемой архитектуре: он разделяет семейство,
критерий и результат оценки, поддерживает MLE, минимальные расстояния,
известную часть параметра и преобразования \cite{reviewDistrMod}.
Его преимущество --- возможность развивать методы, сохраняя общую модель.
Ограничение для PySATL --- зависимость от объектной системы R и
семейства пакетов distr. Расширение distrMod новым оценивателем оставило бы
задачу согласования этих объектов с Python-моделями и вычислениями PySATL.

\subsubsection{OpenTURNS}
OpenTURNS --- наиболее близкая из рассмотренных комплексных
Python-альтернатив. Фабрики распределений позволяют оценивать параметры,
фиксировать известные значения, выбирать решатель и получать
распределение оценок \cite{reviewOTMLE}. Предусмотрены альтернативные
параметризации и пользовательские распределения на Python
\cite{reviewOTParameters,reviewOTPythonDistribution}.
Следовательно, работу с фиксированными параметрами и расширяемость
нельзя считать отсутствующими у этого аналога.

Ограничение связано с ценой принятия его объектной модели.
Оценивание работает с распределениями OpenTURNS. Для уже существующих
семейств PySATL потребуется либо переход на эту модель, либо адаптер,
согласующий параметры, ограничения, вычисления и возвращаемый результат.
Возможность написать такой адаптер есть, но правила согласования
и их проверка остаются частью собственной разработки.
Если бы задача состояла только в подгонке готовых распределений,
использование OpenTURNS было бы самостоятельным решением.

\subsubsection{Специализированные реализации}
В statsmodels можно задать собственный likelihood и моментные условия
GMM \cite{reviewStatsmodelsLikelihood,reviewStatsmodelsGMM}.
Преимущество --- готовые методы и диагностика для этих постановок.
Ограничение --- модель задаётся через соответствующие классы likelihood
или GMM; переход к другому направлению оценивания не образует
единой операции над исходным семейством PySATL.

Lmo в Python и lmomco в R дают развитые средства L-/TL-моментов,
включая модельные значения и подбор параметров
\cite{reviewLmo,reviewLmomco}. Их преимущество --- готовая реализация
специализированных методов, полезная и как независимый эталон.
Ограничение --- моментная специализация: эти пакеты не задают весь
требуемый процесс от произвольного контраста до сопряжённого обновления.
Объединить их с другими средствами можно, но общие правила модели,
параметров и результата потребуется определить отдельно.

\begin{table}[htbp]
\centering\small
\caption{Преимущества и ограничения реализаций для задачи работы}
\label{tab:review-tools}
\begin{tabularx}{\textwidth}{@{}p{0.19\textwidth}XX@{}}
\toprule
Реализация & Преимущества & Остающаяся задача \\
\midrule
SciPy & Распределения, MLE/MM, численные решатели & Объединить методы с разными контрактами \\
fitdistrplus & Несколько методов и законченный процесс подгонки & Связать R-интерфейс с Python и специальными методами \\
distrMod & Разделение семейства, критерия и оценки & Согласовать объектную систему R с моделями PySATL \\
OpenTURNS & Объектное оценивание, фиксация и преобразования параметров & Принять его модель либо поддерживать адаптер \\
statsmodels & Пользовательские likelihood и GMM & Связать разные классы моделей общим процессом \\
Lmo / lmomco & Развитая L-/TL-моментная ветвь & Объединить её с остальными оценивателями \\
\bottomrule
\end{tabularx}
\end{table}

Итак, рынок предоставляет как отдельные алгоритмы, так и развитые
системы подгонки. Недостаток относительно задачи работы --- необходимость
согласовать выбранные методы с одной моделью, параметризациями и
правилами результата. Ни наличие готовой формулы, ни наличие Python API
сами по себе эту задачу не снимают. Конкретные пути её решения
сопоставлены ниже.

\subsection{Обоснование разработки собственного пакета}
Сравнение в разделе~\ref{sec:implementations} показывает, что основная
потребность связана с совместным применением методов и управлением
моделью. Нужна операция, в которой семейство и ограничения задаются
один раз, а оцениватели меняются без повторного описания параметров.
Именно такую проблему выявила работа с фиксированными параметрами
в прототипе.

Например, у гамма-распределения в параметризации shape--rate
среднее равно $m=k/\beta$. При фиксированном $m$ свободным может оставаться
$k$, тогда $\beta=k/m$. В стандартном интерфейсе SciPy используются
shape и scale \cite{reviewScipyGamma}: здесь требуется зависимость
$\mathrm{scale}=m/k$, а не фиксация одного из этих двух параметров.
Её можно реализовать обёрткой или собственным распределением.
Но одинаковая подстановка должна затем работать в MLE, matching,
CDF-расстоянии и каждом повторе бутстрепа. Поэтому нужен общий механизм,
скрывающий эту зависимость от алгоритмов.

\textbf{Расширение SciPy.}
Добавление недостающего метода в SciPy полезно для соответствующей
операции, однако не устраняет различий между контрактами
\texttt{dist.fit}, \texttt{stats.fit} и специализированными средствами.
Для поставленной задачи пришлось бы дополнительно определить общую
цель оценивания, представление свободных параметров, результат и
правила повторного запуска. Это уже самостоятельный интерфейс над
численными процедурами, а не добавление ещё одного метода подгонки.

\textbf{Принятие OpenTURNS или расширение R-пакета.}
OpenTURNS уже решает значительную часть общих задач. Его принятие
означало бы использование ещё одной объектной модели или разработку
адаптера для существующих семейств и прототипов PySATL.
Расширение fitdistrplus либо distrMod, помимо согласования моделей,
сохранило бы необходимость взаимодействия Python и R.
Таким образом, эти варианты переносят часть алгоритмов в готовую систему,
но не устраняют работу над общим представлением задачи в PySATL.

\textbf{Соединение специализированных пакетов.}
Обёртки над statsmodels и Lmo дают доступ к нужным алгоритмам,
но без общего слоя каждая должна согласовывать ограничения, параметры
и результат самостоятельно. При добавлении метода или нового
представления семейства эти правила пришлось бы проверять в нескольких
местах. Собственный пакет позволяет вынести их в одну подготовку задачи
и подключать к ней готовые численные реализации.

Поэтому объём собственной разработки определяется оставшимися общими
задачами: отделить свободные параметры от фиксированных, предоставить
оценивателю нужные характеристики модели, согласовать результаты
разных методов и повторное оценивание. Результат должен сохранять
способ получения оценки и сведения о численных неудачах. Решатели SciPy и применимые
специализированные вычисления при этом переиспользуются.
Обоснован именно собственный пакет согласованного оценивания,
а не повторная реализация всех возможностей рассмотренных библиотек.

Прототипы уже дают основу MLE и эмпирических распределений
\cite{reviewPrototype}. Их развитие позволяет проверить предложенный
общий механизм на существующих и добавляемых методах. Практическая
ценность пакета должна проявиться в сквозном сценарии: один раз задать
семейство и фиксированные параметры, сменить оцениватель, выполнить
бутстреп и получить сопоставимые результаты с понятной диагностикой.

\subsection{Анализ и выводы}
Таблица~\ref{tab:review-tools} показывает, что готовые реализации покрывают
значительную часть выбранных методов. Для сохранения графа характеристик,
параметризаций и ограничений PySATL пришлось бы согласовывать в адаптере
две модели распределений, преобразования параметров и диагностику.
Для данной задачи этот путь оценивается как более сложный в разработке
и менее эффективный в сопровождении, чем реализация в Core.
Поэтому выбирается общий слой в Core, использующий его модель напрямую
и готовые численные процедуры.

Из различий методов в разделе~\ref{sec:methods} следует необходимость
разделить статистическую постановку и способ вычисления оценки.
Минимизация критерия, решение уравнений и сопряжённое обновление требуют
разных вычислительных путей. Общий интерфейс должен предоставлять
каждому методу нужные характеристики модели и только свободные параметры.
Подстановка фиксированных значений остаётся внутри модели; проверка
допустимости и представление результата используются всеми методами.
Такая организация позволяет повторно применять оцениватель в бутстрепе.

Для реализации выбираются Python и существующая модель PySATL Core,
а для численной оптимизации и решения уравнений --- процедуры SciPy.
fitdistrplus, OpenTURNS, Lmo и lmomco полезны как независимые реализации
для проверки соответствующих методов. Сравнение должно учитывать
одинаковые параметризации, ограничения, определения статистик и настройки
решателей.
\label{sec:relatedworks-end}
